\documentclass[aspectratio=169,11pt,english]{beamer}
\input{../../../_preambulo_beamer.tex}
\input{../../../_comandos_beamer.tex}
\title[L05 · 40 minutes]{Set Theory and Sample Spaces}
\begin{document}
\begin{frame}{L05 · One dispatch, three features}
\framesubtitle{00–02 min · Case}
A dispatch records its \textbf{shift} (Day/Night), \textbf{service} (Standard/Express), and \textbf{status} (On time/Delayed).

\medskip
\textbf{Question:} which dispatches can be at night or delayed?

\medskip
We use eight possible types and frequencies from 480 \textbf{synthetic} dispatches. They do not represent a real company.
\end{frame}
\begin{frame}{One outcome is a tuple; $S$ collects the possibilities}
\framesubtitle{02–07 min · Model}
\[
\omega=(\text{shift},\text{service},\text{status}),\qquad S=T\times V\times E.
\]
$T=\{\text{Day},\text{Night}\}$, $V=\{\text{Standard},\text{Express}\}$, and $E=\{\text{On time},\text{Delayed}\}$.
\[
|S|=2\cdot2\cdot2=8.
\]
\textbf{Check:} what are the two Night, Express tuples?

\medskip
\textbf{Answer:} (Night, Express, On time) and (Night, Express, Delayed).
\end{frame}
\begin{frame}{Notebook · Enumerate the eight tuples}
\framesubtitle{07–10 min · Cell 1}
\texttt{product(SHIFTS, SERVICES, STATUSES)} builds $S$ in category order.

\medskip
\textbf{Run cell 1.} Check that it shows eight distinct tuples. The first is (Day, Standard, On time); the last is (Night, Express, Delayed).

\medskip
A tuple identifies an \textbf{outcome type}, not an individual dispatch.
\end{frame}
\begin{frame}{Eight possibilities; 480 records}
\framesubtitle{10–14 min · Cell 2}
Cell 2 assigns synthetic frequencies to the eight tuples. The extended script generated them with seed 42.
\[
\sum_{\omega\in S}n(\omega)=480,\qquad |S|=8.
\]
For example, (Day, Standard, On time) occurs 196 times.

\medskip
\textbf{Check:} if another dispatch has that type, does $|S|$ change?

\medskip
\textbf{Answer:} no; $n(\omega)$ increases, while the number of possibilities stays eight.
\end{frame}
\begin{frame}{Events select outcome types}
\framesubtitle{14–18 min · Cell 3}
\[
A=\{\omega\in S:\text{shift = Night}\},\qquad
B=\{\omega\in S:\text{status = Delayed}\}.
\]
$A$ and $B$ are subsets of $S$, each containing four types.

\medskip
\textbf{Before running:} does (Night, Express, Delayed) belong to both?

\medskip
Yes. The same tuple can satisfy two conditions.
\end{frame}
\begin{frame}{Intersection, union, and complement}
\framesubtitle{18–23 min · Prediction}
\[
\begin{array}{ll}
A\cap B &: \text{Night and Delayed}\quad |A\cap B|=2,\\[3pt]
A\cup B &: \text{Night or Delayed}\quad |A\cup B|=6,\\[3pt]
A^c=S\setminus A &: \text{Day}\quad |A^c|=4.
\end{array}
\]
The \textbf{or} is inclusive: a tuple satisfying both conditions appears once in $A\cup B$.

\medskip
\textbf{Check:} $|A\cup B|=4+4-2=6$.
\end{frame}
\begin{frame}{Notebook · Possible types and observed records}
\framesubtitle{23–26 min · Cell 3}
\centering
\begin{tabular}{lrr}
\toprule Event & Possible types & Synthetic records\\ \midrule
$A$ & 4 & 172\\
$B$ & 4 & 86\\
$A\cap B$ & 2 & 45\\
$A\cup B$ & 6 & 213\\
$A^c$ & 4 & 308\\ \bottomrule
\end{tabular}
\par\medskip\raggedright
\textbf{Run cell 3.} Why is 45 not the cardinality of $A\cap B$?

\medskip
Because 45 counts observed dispatches; the intersection contains two types.
\end{frame}
\begin{frame}{Disjointness and exhaustiveness ask different questions}
\framesubtitle{26–30 min · Definitions}
A pair is \textbf{disjoint} if $X\cap Y=\varnothing$. It is \textbf{exhaustive} if $X\cup Y=S$.

\medskip
For $A$ (Night) and $B$ (Delayed):
\[
|A\cap B|=2,\qquad |A\cup B|=6.
\]
They are neither disjoint nor exhaustive.

\medskip
\textbf{Question:} which two types are outside $A\cup B$?

\medskip
The Day and On time types, one for each service.
\end{frame}
\begin{frame}{Notebook · A pair that partitions $S$}
\framesubtitle{30–33 min · Cell 4}
Let $C$ = Standard and $D$ = Express.
\[
C\cap D=\varnothing,\qquad C\cup D=S.
\]
\textbf{Run cell 4.} Compare $(A,B)$ with $(C,D)$.

\medskip
\centering
\begin{tabular}{lcc}
\toprule Pair & Disjoint & Exhaustive\\ \midrule
$A,B$ & No & No\\
$C,D$ & Yes & Yes\\ \bottomrule
\end{tabular}
\end{frame}
\begin{frame}{Integrate · A new question}
\framesubtitle{33–37 min · Cell 5}
Define $F$ = Express and Delayed.

\medskip
\textbf{Before running:} predict $|F|$, the number of synthetic records in $F$, $|F\cap A^c|$, and $|F\cup A^c|$.

\medskip
Hint: $A^c$ is the set of Day outcomes.

\medskip
Run cell 5 and explain each number in its unit: possible type or observed record.
\end{frame}
\begin{frame}{Check the answer}
\framesubtitle{37–39 min · Interpretation}
\[
|F|=2,\qquad n(F)=40,\qquad
|F\cap A^c|=1,\qquad |F\cup A^c|=5.
\]
One tuple is Express, Delayed, and Day. The other is Express, Delayed, and Night.

\medskip
\textbf{Exit question:} why do $|F|$ and $n(F)$ answer different questions?
\end{frame}
\begin{frame}{Continue studying}
\framesubtitle{39–40 min · Close}
The extended notebook develops the simulation of 480 dispatches, figures, and further exercises. The three scripts show how records are generated and analyzed.

\medskip
\textbf{Bridge to L06:} as the sample space grows, listing every outcome becomes impractical. We will need counting rules.
\end{frame}
\end{document}
